\section{Интерактивный сервис оценки}
\label{sec:service}

\subsection{Назначение}

Помимо опубликованного набора данных и обученных моделей, практическая ценность работы
включает интерактивный сервис, позволяющий применять разработанные оценщики к новым
RAG-ответам без необходимости обращаться к обучающему конвейеру напрямую. Сервис
реализован на Streamlit (модуль \texttt{new\_streamlit.py}) и запускается командой
\texttt{streamlit run new\_streamlit.py}. Применение модели может выполняться как на CPU, так и
на одном GPU, без специальной дополнительной настройки.

Роль сервиса в контексте диссертации прикладная: он не является центральным научным
результатом работы, но демонстрирует, что обученные модели можно использовать не только
в экспериментальном конвейере, но и в режиме анализа новых примеров.

\subsection{Архитектура}

Логически сервис состоит из трех блоков.

\begin{itemize}
  \item \textbf{Реестр моделей.}
  В сервисе задан список доступных конфигураций, соответствующих обученным моделям,
  предназначенным для практического использования. На момент написания работы доступны
  три базовые конфигурации:
    \begin{itemize}
      \item \texttt{DeBERTa-v3-large (len=512)} - контрольная базовая модель;
      \item \texttt{ModernBERT (len=1024)} - основная русскоязычная конфигурация;
      \item \texttt{BAAI-bge-m3 (len=1024)} - многоязычная модель для смешанных корпусов.
    \end{itemize}
  Веса хранятся локально и повторно используются средствами Streamlit, что позволяет не
  перегружать память повторной инициализацией модели при переключении между сессиями.

  \item \textbf{Модуль применения модели.}
  Внутри сервиса реализован легковесный модуль загрузки сохранённых весов, совместимый с
  обучающей реализацией модели-оценщика. На вход он принимает тройку
  <<вопрос -- контекст -- ответ>> в формате, близком к структуре RAGBench-RU, а на
  выходе формирует вероятности и агрегированные значения TRACe-метрик.

  \item \textbf{Пользовательский интерфейс.}
  Интерфейс поддерживает два основных режима:
    \begin{enumerate}
      \item \emph{Одиночный режим} - оценка одной тройки
      (вопрос, контекст, ответ) с визуализацией результата;
      \item \emph{Пакетный режим} - оценка CSV/JSONL-файла с последующей выгрузкой
      рассчитанных оценок.
    \end{enumerate}
\end{itemize}

\subsection{Основные возможности}

\begin{itemize}
  \item \textbf{Одиночная оценка примера.}
  Пользователь может ввести вопрос, контекст и ответ вручную и получить три основных
  значения TRACe-метрики, а также вероятности на уровне токенов. Такой режим полезен для разбора
  конкретных ошибок и качественного анализа поведения модели.

  \item \textbf{Пакетная обработка.}
  Сервис поддерживает загрузку таблиц с заранее подготовленными примерами и возвращает
  результаты в виде CSV с рассчитанными значениями
  \texttt{relevance\_score}, \texttt{utilization\_score} и
  \texttt{adherence\_score}. Это делает сервис удобным для автономного анализа небольших
  коллекций ответов.

  \item \textbf{Настройка порогов.}
  Пороги $\tau_R$, $\tau_U$, $\tau_A$ можно изменять интерактивно. Такая возможность
  важна не только с прикладной точки зрения, но и как средство исследования: она
  позволяет наблюдать, как меняется визуальная разметка при переходе от более
  консервативного к более чувствительному режиму.

  \item \textbf{Визуализация разметки на уровне токенов.}
  В одиночном режиме документ и ответ отображаются с цветовой подсветкой, отражающей
  предсказания релевантности, использованности и верности. Это помогает интерпретировать модель
  не только через итоговые оценки, но и через локальные решения на уровне текста.

  \item \textbf{Совместимость с предобученными моделями и локально сохранёнными весами.}
  Сервис загружает базовые модели и токенизаторы через интерфейс Hugging Face
  Transformers, а дообученные веса TRACe-оценщика считывает из локально сохраненных
  файлов. Поэтому он может использоваться как интерфейс к уже полученным
  экспериментальным результатам при наличии соответствующих локальных артефактов
  модели и файла с подобранными порогами.


\end{itemize}

На рисунке~\ref{fig:service_example} приведен пример работы сервиса на одном из
примеров набора данных RAGBench-RU. В интерфейсе выбраны модель, пример из
набора данных и пороги классификации, после чего сервис рассчитывает значения
TRACe-метрик и сопоставляет их с эталонной разметкой. При выбранных порогах
оценки релевантности, использованности и верности близки к эталонным значениям.
Кроме того, таблица вероятностей на уровне токенов показывает, какие фрагменты
текста модель считает наиболее значимыми для каждой метки: более насыщенный цвет
ячейки соответствует большей вероятности.

\begin{figure}[H]
  \centering
  \includegraphics[width=\textwidth]{service_example.png}
  \caption{Пример работы интерактивного сервиса на примере из RAGBench-RU:
  сравнение рассчитанных TRACe-метрик с эталонной разметкой и визуализация
  вероятностей на уровне токенов.}
  \label{fig:service_example}
\end{figure}

\subsection{Ограничения}

Сервис сознательно не пытается решать все задачи, связанные с жизненным циклом
моделей-оценщиков.

\begin{itemize}
  \item Он ориентирован на одиночную и пакетную автономную оценку, а не на полноценное
  производственное развертывание.
  \item Он не содержит встроенного механизма дообучения на пользовательских данных;
  для таких сценариев следует использовать основной обучающий конвейер проекта.
  \item В работе не проводились формальные измерения производительности сервиса
  (задержка, пропускная способность, устойчивость при высокой нагрузке), поэтому
  приведенная реализация рассматривается как прикладной интерфейс, а не как готовый
  промышленный компонент.
  \item Интеграция через REST-API, webhook или CI-контуры в текущей версии не
  реализована и рассматривается как направление дальнейшего развития.
\end{itemize}

Таким образом, Streamlit-сервис выступает в работе как удобный мост между
исследовательским результатом и его практическим использованием: он показывает, что
обученные оценщики можно применять к новым данным уже в текущем виде, но не подменяет
собой отдельную задачу сервисного развертывания.
